% =============================================================================
%  s03_conclusiones.tex — Sección 3: Conclusiones y recomendaciones
% =============================================================================
\section{Conclusiones}

% ── Frame: Solo texto — mensaje clave, cita o puente entre secciones ─────────
%    USO: para declaraciones fuertes, preguntas retóricas o transiciones.
%    Copia y adapta. No lleva título de frame ni elementos gráficos.
\begin{frame}[plain]
  \begin{tikzpicture}[remember picture, overlay]
    \fill[ColorPrincipal]
      (current page.north west) rectangle
      ([xshift=0.35cm] current page.south west);
    \fill[ColorAcento]
      ([xshift=0.32cm] current page.north west) rectangle
      ([xshift=0.35cm] current page.south west);
  \end{tikzpicture}
  \hspace{0.8cm}%
  \begin{minipage}{0.88\textwidth}
    \vspace{2.2cm}
    {\color{ColorAcento}\rule{0.5\textwidth}{1.2pt}}\\[0.5cm]
    {\large\bfseries\color{ColorPrincipal}
      ¿Es posible la Movilidad Sustentable en la CDMX y su ZMVM?\\
      ¿Por qué no hay concordancia entre la teoría y la realidad?}\\[0.6cm]
    {\normalsize\color{black}
      La \cdmx  y la \zmvm  no está perdida, pero la movilidad ha sido reducida a
      propuestas de manual y se ha segregado a los sectores populares a ciudadanos de
      segunda o tercera clase \cite{rendon2009ciudadanos}. No se requiere soluciones vistosas
       o políticas, si no socialmente incluyentes y con una contextualización del entorno adecuada.
      }\\[0.5cm]
    {\color{ColorAcento}\rule{0.3\textwidth}{0.8pt}}
    \vfill
  \end{minipage}
\end{frame}


% ── Frame: Síntesis de Hallazgos (I) ──────────────────────────────────────────
\begin{frame}{Síntesis de Hallazgos: Eficiencia y Equidad Lograda}
  \begin{columns}[T]
    \begin{column}{0.90\textwidth}
      \begin{block}{Lo que funciona}
        \begin{itemize}
          \item \textbf{Superioridad del Transporte Masivo:} La ingeniería vial demuestra que el transporte masivo mueve hasta 
              \textbf{43,000 pers/hr} frente a las 2,000 del auto, validando los principios de Transportar y Densificar \cite{ITDP_Espacio}.
          \item \textbf{Tecnología como Justicia Social:} El Cablebús ha reducido tiempos de traslado en un \textbf{45\%} 
                en periferias, demostrando que la innovación tecnológica puede mitigar la segregación socioespacial \cite{SEDATU2023}.
          \item \textbf{Sustento Normativo Sólido:} La ENAMOV 2023-2042 y el Artículo 4° brindan el blindaje legal necesario para transitar 
                de una planeación de tráfico a una de \textbf{derecho a la movilidad} \cite{SEDATU2023,ITDP_Espacio}.
        \end{itemize}
      \end{block}
    \end{column}
  \end{columns}
\end{frame}

% ── Frame: Síntesis de Hallazgos (II) ─────────────────────────────────────────
\begin{frame}{Síntesis de Hallazgos: Retos y Barreras Estructurales}
  \begin{columns}[T]
    \begin{column}{0.98\textwidth}
      \begin{alertblock}{Los retos pendientes}
        \begin{itemize}
          \item \textbf{Descontrol de la Motorización:} El parque vehicular crece al \textbf{5.3\% anual} 
                (3.5 veces más que la población), lo que anula las ganancias en eficiencia vial si no se 
                reduce el uso del auto \cite{SEDATU2023}.
          \item \textbf{Agotamiento del Modelo Radial:} La centralización operativa genera altos 
                \textbf{Coeficientes de Fricción ($C_f$)}, saturando hubs como Pantitlán y dejando al 
                  4to contorno con coberturas marginales ($<12\%$) \cite{Garibaldi2024_R02}.
          \item \textbf{Costo de la Ineficiencia:} Las externalidades negativas (congestión y siniestros) 
                consumen el \textbf{4\% del PIB}, impulsadas por un uso ineficiente del espacio vial 
                (68.1\% de ocupación individual) \cite{ramirez2025analisis}.
        \end{itemize}
      \end{alertblock}
    \end{column}
  \end{columns}
\end{frame}

% ── Frame: Cierre ────────────────────────────────────────────────────────────
\begin{frame}[plain]
  \begin{center}
    \vfill
    {\color{ColorAcento}\rule{0.6\textwidth}{1.5pt}}\\[0.5cm]
    {\Large\bfseries\color{ColorPrincipal} Gracias}\\[0.3cm]
    {\normalsize \Autor}\\[0.1cm]
    {\small \MateriaCorta}\\[0.1cm]
    {\small \Semestre}\\[0.5cm]
    {\color{ColorAcento}\rule{0.6\textwidth}{1.5pt}}
    \vfill
  \end{center}
\end{frame}
